\chapter{Sub-exponential random variables}
\label{chap:pre_subexp}

This chapter develops the Bernstein-type concentration used by the $\ell_2$-norm estimation
of Chapter~\ref{chap:norm} and by the unit-ball identification algorithm
(Lemma~\ref{lem:ball_identification}): a notion of sub-exponential random variable with two
explicit parameters $(V, b)$ (Appendix~D.1 of the paper), its stability under scaling and
independent sums, the corresponding tail bound with the explicit constant $c = 1/2$, and the
sub-exponential parameters of the centered squares of Gaussian variables (the paper's
Lemma~1 in Appendix~A, with its complete-the-square proof), of chi-square variables and of
Rademacher linear forms $\ip{\eps}{\theta}$. For the last one the paper invokes the
Hanson--Wright inequality with unspecified constants; we replace it by a direct argument valid
for every sub-Gaussian variable $Z$ with $\E Z^2$ equal to its sub-Gaussian constant (a Gaussian
trick for positive $\lambda$, a fourth-moment bound for negative $\lambda$, the fourth moment
being itself controlled by the Gaussian trick), which gives explicit parameters
$(16\norm\theta^4, 4\norm\theta^2)$. All constants are explicit and tracked. References:
\cite[Section~2.7]{vershynin2018high}, \cite[Section~2.4]{boucheron2013concentration}.

\textbf{What Mathlib provides.} The sub-Gaussian structure
\texttt{ProbabilityTheory.HasSubgaussianMGF X c μ} (\texttt{Mathlib/Probability/Moments/SubGaussian.lean})
with \texttt{measure\_ge\_le}, \texttt{neg}, \texttt{add\_of\_indepFun},
\texttt{sum\_of\_iIndepFun}, \texttt{of\_map}, \texttt{congr\_identDistrib}, Hoeffding's lemma
\texttt{hasSubgaussianMGF\_of\_mem\_Icc\_of\_integral\_eq\_zero}, and LML's
\texttt{measure\_le\_le}. Moment generating functions: \texttt{ProbabilityTheory.mgf},
\texttt{iIndepFun.mgf\_sum}, \texttt{iIndepFun.integrable\_exp\_mul\_sum}, the Chernoff bound
\texttt{measure\_ge\_le\_exp\_mul\_mgf} (\texttt{Mathlib/Probability/Moments/Basic.lean}),
the derivative of the MGF \texttt{hasDerivAt\_mgf} (\texttt{MGFAnalytic.lean}), integrability
of all moments from the integrability of $e^{\lambda X}$ near $\lambda = 0$
(\texttt{integrable\_pow\_of\_mem\_interior\_integrableExpSet}, \texttt{IntegrableExpMul.lean}),
\texttt{mgf\_gaussianReal}. The Gaussian integral \texttt{integral\_gaussian}
($\int e^{-bx^2}dx = \sqrt{\pi/b}$,
\texttt{Mathlib/Analysis/SpecialFunctions/Gaussian/GaussianIntegral.lean}); the real logarithm
series bound \texttt{Real.abs\_log\_sub\_add\_sum\_range\_le}
(\texttt{Mathlib/Analysis/SpecialFunctions/Log/Deriv.lean}); \texttt{Real.add\_one\_le\_exp} and
\texttt{Real.quadratic\_le\_exp\_of\_nonneg}; independence of the coordinates of a product
measure \texttt{ProbabilityTheory.iIndepFun\_pi}; the standard Gaussian vector
\texttt{ProbabilityTheory.stdGaussian} as the image of a product of \texttt{gaussianReal 0 1}
(\texttt{map\_pi\_eq\_stdGaussian}).

\textbf{What is missing.} The sub-exponential structure itself and all its properties
(Section~\ref{sec:pre_subexp_def}), the MGF of the square of a real Gaussian
(Lemma~\ref{lem:gaussian_square_mgf}), and the sub-exponential bounds of
Sections~\ref{sec:pre_subexp_gaussian} and~\ref{sec:pre_subexp_rademacher}.

\textbf{Lean remark.} The definition below mirrors \texttt{HasSubgaussianMGF}: a structure
\texttt{HasSubexponentialMGF (X : Ω → ℝ) (V b : ℝ) (μ : Measure Ω) : Prop} with fields
\texttt{integrable\_exp\_mul : ∀ t : ℝ, b * |t| ≤ 1 → Integrable (fun ω ↦ exp (t * X ω)) μ} and
\texttt{mgf\_le : ∀ t : ℝ, b * |t| ≤ 1 → mgf X μ t ≤ exp (V * t\^{}2 / 2)}
(\texttt{COLT83/Mathlib/Probability/Moments/SubExponential.lean}). The range of $t$ is
written as $b\abs t\le1$ rather than $\abs t\le1/b$ so that \texttt{b = 0} imposes no
restriction on $t$ (the hypothesis $0\cdot\abs t\le1$ is always true):
\texttt{HasSubexponentialMGF X V 0 μ} is then equivalent to \texttt{HasSubgaussianMGF X V μ}
(same two fields, quantified over all \texttt{t}), and Mathlib's sub-Gaussian API is the
$b = 0$ case of the statements below (Lemma~\ref{lem:subexp_of_subgaussian}). Unlike
Mathlib's sub-Gaussian constant, the parameters are real numbers (not \texttt{ℝ≥0}), which
avoids coercions in the arithmetic of Chapter~\ref{chap:norm}; the lemmas that need $V > 0$
or $b > 0$ assume it explicitly. A kernel version, as for \texttt{Kernel.HasSubgaussianMGF}, is
not needed.

\section{Definition and basic properties}
\label{sec:pre_subexp_def}

\begin{definition}[Sub-exponential moment generating function]\label{def:subexp}
\lean{ProbabilityTheory.HasSubexponentialMGF}\leanok
Let $X : \Omega \to \R$ be a random variable on $(\Omega,\Prob)$ and $V, b \in \R$. We say
that $X$ \emph{has a sub-exponential MGF with parameters $(V,b)$}, written
$\mathrm{HasSubexponentialMGF}(X, V, b)$, if for every $\lambda \in \R$ with
$b\abs{\lambda} \le 1$ (i.e.\ $\abs\lambda\le1/b$ when $b>0$, and every $\lambda\in\R$ when
$b \le 0$), the variable $e^{\lambda X}$ is integrable and
\[
  \E e^{\lambda X} \le \exp\Big(\frac{\lambda^2V}{2}\Big).
\]
The case $b = 0$ is exactly Mathlib's \texttt{HasSubgaussianMGF X V}
(Lemma~\ref{lem:subexp_of_subgaussian}). No centering is assumed; it is a consequence
(Lemma~\ref{lem:ps_subexp_integral_zero}). In practice $V \ge 0$ and $b \ge 0$.
\end{definition}

\begin{lemma}[Sub-exponential variables are integrable and centered]\label{lem:ps_subexp_integral_zero}
\lean{ProbabilityTheory.HasSubexponentialMGF.zero_mem_interior_integrableExpSet, ProbabilityTheory.HasSubexponentialMGF.integrable, ProbabilityTheory.HasSubexponentialMGF.memLp, ProbabilityTheory.HasSubexponentialMGF.integrable_pow, ProbabilityTheory.HasSubexponentialMGF.integral_eq_zero}\leanok
\uses{def:subexp}
If $\mathrm{HasSubexponentialMGF}(X,V,b)$ then $X$ is integrable, all its moments are finite,
and, when $\Prob$ is a probability measure, $\E X = 0$.
\end{lemma}

\begin{proof}
\leanok
\uses{def:subexp}
The set $\{\lambda : b\abs\lambda \le 1\}$ is a neighbourhood of $0$ (for every $b$, since
$\lambda \mapsto b\abs\lambda$ is continuous and vanishes at $0$), so $0$ lies in the interior
of the set of $\lambda$ such that $e^{\lambda X}$ is integrable. Mathlib then provides the
integrability of $\abs X^p$ for every $p$
(\texttt{integrable\_pow\_of\_mem\_interior\_integrableExpSet}) and the differentiability of
$\lambda \mapsto \E e^{\lambda X}$ at $0$ with derivative $\E X$ (\texttt{hasDerivAt\_mgf}).
Centering: the function $\varphi(\lambda) := \E e^{\lambda X} - e^{\lambda^2V/2}$ is
nonpositive on a neighbourhood of $0$ and vanishes at $0$ (as $\E e^{0} = 1$ for a probability
measure), so it has a local maximum at $0$; it is differentiable there with derivative
$\E X - 0$, which must vanish (Fermat's rule, Mathlib \texttt{IsLocalMax.hasDerivAt\_eq\_zero}).
Hence $\E X = 0$.
\end{proof}

\begin{lemma}[Monotonicity in the parameters]\label{lem:ps_subexp_mono}
\lean{ProbabilityTheory.HasSubexponentialMGF.mono}\leanok
\uses{def:subexp}
If $\mathrm{HasSubexponentialMGF}(X,V,b)$, $V \le V'$ and $b \le b'$, then
$\mathrm{HasSubexponentialMGF}(X,V',b')$.
\end{lemma}

\begin{proof}
\leanok
\uses{def:subexp}
$b'\abs\lambda \le 1$ implies $b\abs\lambda \le b'\abs\lambda \le 1$, and
$e^{\lambda^2V/2} \le e^{\lambda^2V'/2}$.
\end{proof}

\begin{lemma}[Sub-Gaussian is sub-exponential with $b = 0$]\label{lem:subexp_of_subgaussian}
\lean{ProbabilityTheory.HasSubexponentialMGF.of_hasSubgaussianMGF, ProbabilityTheory.HasSubexponentialMGF.hasSubgaussianMGF}\leanok
\uses{def:subexp, lem:ps_subexp_mono}
$X$ has a sub-Gaussian MGF with constant $c \ge 0$ (Mathlib
\texttt{HasSubgaussianMGF X c}: $e^{\lambda X}$ integrable and $\E e^{\lambda X} \le e^{c\lambda^2/2}$
for all $\lambda\in\R$) if and only if $\mathrm{HasSubexponentialMGF}(X, c, 0)$. Consequently,
if $X$ has a sub-Gaussian MGF with constant $c$, then $\mathrm{HasSubexponentialMGF}(X, c, b)$
for every $b \in \R$.
\end{lemma}

\begin{proof}
\leanok
\uses{def:subexp, lem:ps_subexp_mono}
For $b = 0$ the condition $b\abs\lambda\le1$ holds for every $\lambda\in\R$, so the two
requirements of Definition~\ref{def:subexp} are literally the two fields of
\texttt{HasSubgaussianMGF}. The second claim is the same argument: the condition
$b\abs\lambda \le 1$ only restricts the set of $\lambda$ for which the two fields are required.
\end{proof}

\begin{lemma}[Scaling]\label{lem:subexp_const_mul}
\lean{ProbabilityTheory.HasSubexponentialMGF.const_mul, ProbabilityTheory.HasSubexponentialMGF.neg}\leanok
\uses{def:subexp}
If $\mathrm{HasSubexponentialMGF}(X,V,b)$ and $a \in \R$, then
$\mathrm{HasSubexponentialMGF}(aX,\ a^2V,\ \abs{a}\,b)$. In particular ($a = -1$),
$\mathrm{HasSubexponentialMGF}(-X, V, b)$.
\end{lemma}

\begin{proof}
\leanok
\uses{def:subexp}
Let $(\abs ab)\abs\lambda \le 1$, i.e.\ $b\abs{\lambda a} \le 1$; then $e^{\lambda(aX)} = e^{(\lambda a)X}$
is integrable and $\E e^{(\lambda a)X} \le e^{(\lambda a)^2V/2} = e^{\lambda^2(a^2V)/2}$
(for $a = 0$ this reads $\E e^{0} = 1 \le 1$).
\end{proof}

\begin{lemma}[Transport]\label{lem:ps_subexp_map}
\lean{ProbabilityTheory.HasSubexponentialMGF.congr, ProbabilityTheory.HasSubexponentialMGF.of_map, ProbabilityTheory.HasSubexponentialMGF.map_iff, ProbabilityTheory.HasSubexponentialMGF.id_map_iff, ProbabilityTheory.HasSubexponentialMGF.congr_identDistrib}\leanok
\uses{def:subexp}
The property $\mathrm{HasSubexponentialMGF}(X,V,b)$ depends only on the law of $X$: if
$X = Y$ almost surely, or if $X$ and $Y$ have the same law (on possibly different probability
spaces), then $X$ has parameters $(V,b)$ if and only if $Y$ does. If $Y : \Omega' \to \Omega$ is
measurable and $X : \Omega \to \R$ has parameters $(V,b)$ under the image measure
$\Prob' \circ Y^{-1}$, then $X \circ Y$ has parameters $(V,b)$ under $\Prob'$.
\end{lemma}

\begin{proof}
\leanok
\uses{def:subexp}
Both fields of Definition~\ref{def:subexp} are expressed through integrals of functions of
$X$, which are unchanged by these operations (change of variables
\texttt{integrable\_map\_measure}, \texttt{mgf\_map}). This is the same proof as for Mathlib's
\texttt{HasSubgaussianMGF.of\_map} and \texttt{congr\_identDistrib}.
\end{proof}

\begin{lemma}[Bernstein-type tail bound]\label{lem:subexp_tail}
\lean{ProbabilityTheory.HasSubexponentialMGF.measure_ge_le_exp_neg_sq, ProbabilityTheory.HasSubexponentialMGF.measure_ge_le_exp_neg_div, ProbabilityTheory.HasSubexponentialMGF.measure_ge_le, ProbabilityTheory.HasSubexponentialMGF.measure_le_le, ProbabilityTheory.HasSubexponentialMGF.measure_abs_ge_le, ProbabilityTheory.HasSubexponentialMGF.measure_abs_gt_le}\leanok
\uses{def:subexp}
Let $\mathrm{HasSubexponentialMGF}(X,V,b)$ and $t \ge 0$. Then
\[
  \Prob(X \ge t) \le
  \begin{cases}
    \exp\big(-\frac{t^2}{2V}\big) & \text{if } tb \le V \text{ and } V > 0,\\[2pt]
    \exp\big(-\frac{t}{2b}\big) & \text{if } tb \ge V \text{ and } b > 0 ,
  \end{cases}
\]
and consequently, when $V > 0$ and $b > 0$,
\[
  \Prob(X \ge t) \le \exp\Big(-\min\Big(\frac{t^2}{2V}, \frac{t}{2b}\Big)\Big), \quad
  \Prob(X \le -t) \le \exp\Big(-\min\Big(\frac{t^2}{2V}, \frac{t}{2b}\Big)\Big), \quad
  \Prob(\abs X \ge t) \le 2\exp\Big(-\min\Big(\frac{t^2}{2V}, \frac{t}{2b}\Big)\Big).
\]
If $b = 0$ (the sub-Gaussian case) and $V > 0$, the first bound applies for every $t \ge 0$
(the condition $tb \le V$ is automatic): this is Mathlib
\texttt{HasSubgaussianMGF.measure\_ge\_le} (and LML \texttt{measure\_le\_le} for the lower tail),
i.e.\ the $\min$ form with the convention $t/(2b) = +\infty$.
\end{lemma}

\begin{proof}
\leanok
\uses{def:subexp, lem:subexp_const_mul}
For $\lambda \ge 0$ with $b\lambda \le 1$, Markov's inequality applied to the nonnegative
integrable variable $e^{\lambda X}$ gives (Chernoff, Mathlib
\texttt{measure\_ge\_le\_exp\_mul\_mgf})
\[
  \Prob(X \ge t) = \Prob(e^{\lambda X} \ge e^{\lambda t}) \le e^{-\lambda t}\,\E e^{\lambda X}
  \le \exp\Big(-\lambda t + \frac{\lambda^2V}{2}\Big).
\]
\emph{Case $tb \le V$, $V > 0$:} take $\lambda := t/V \ge 0$, which satisfies $b\lambda \le 1$;
the exponent is $-t^2/V + t^2/(2V) = -t^2/(2V)$. \emph{Case $tb \ge V$, $b > 0$:} take
$\lambda := 1/b$; the exponent is $-t/b + V/(2b^2) \le -t/b + t/(2b) = -t/(2b)$ since
$V \le tb$. For the $\min$ form it is not necessary to identify which term realizes the
minimum: in each case the bound obtained is $\exp(-u)$ for one of the two terms $u$ of the
minimum, and $\exp(-u) \le \exp(-\min(\cdot,\cdot))$. The lower tail is the upper tail of
$-X$, which has the same parameters (Lemma~\ref{lem:subexp_const_mul} with $a = -1$); the
two-sided bound is a union bound over $\{X \ge t\} \cup \{X \le -t\}$.
\end{proof}

\begin{remark}[The constant $c$ of the paper]\label{rem:ps_tail_constant}
The paper's Lemma~2 (Appendix~A) states this bound with an unspecified absolute constant $c$
in the exponent; the computation above, which is the one given there, yields $c = 1/2$. This
value is used in Chapter~\ref{chap:norm} to make all sample-complexity constants explicit.
\end{remark}

\begin{lemma}[Independent sums]\label{lem:subexp_sum}
\lean{ProbabilityTheory.HasSubexponentialMGF.add_of_indepFun, ProbabilityTheory.HasSubexponentialMGF.sum_of_iIndepFun, ProbabilityTheory.HasSubexponentialMGF.fun_sum_of_iIndepFun}\leanok
\uses{def:subexp}
Let $X_1,\dots,X_K$ be independent with $\mathrm{HasSubexponentialMGF}(X_i,V_i,b)$ for every
$i \le K$ (the same $b$). Then $\mathrm{HasSubexponentialMGF}\big(\sum_{i\le K}X_i,\ \sum_{i\le K}V_i,\ b\big)$.
\end{lemma}

\begin{proof}
\leanok
\uses{def:subexp}
Let $b\abs\lambda \le 1$. Each $e^{\lambda X_i}$ is integrable and they are independent, so
$e^{\lambda\sum_iX_i} = \prod_ie^{\lambda X_i}$ is integrable (Mathlib
\texttt{iIndepFun.integrable\_exp\_mul\_sum}) and
$\E e^{\lambda\sum_iX_i} = \prod_i\E e^{\lambda X_i}$ (\texttt{iIndepFun.mgf\_sum})
$\le \prod_ie^{\lambda^2V_i/2} = e^{\lambda^2\sum_iV_i/2}$.
\end{proof}

\begin{corollary}[Tail of an average]\label{cor:subexp_average_tail}
\lean{ProbabilityTheory.HasSubexponentialMGF.average_of_iIndepFun, ProbabilityTheory.HasSubexponentialMGF.measure_abs_average_ge_le, ProbabilityTheory.HasSubexponentialMGF.measure_abs_average_ge_le_of_forall}\leanok
\uses{def:subexp}
Let $K \ge 1$ and $X_1,\dots,X_K$ be independent with $\mathrm{HasSubexponentialMGF}(X_i,V_i,b)$,
and let $\bar V := \frac1K\sum_{i\le K}V_i$. Then
$\mathrm{HasSubexponentialMGF}\big(\frac1K\sum_{i\le K}X_i,\ \bar V/K,\ b/K\big)$ and, if
$\bar V > 0$ and $b > 0$, for every $t \ge 0$,
\[
  \Prob\Big(\Big|\frac1K\sum_{i\le K}X_i\Big| \ge t\Big)
  \le 2\exp\Big(-K\min\Big(\frac{t^2}{2\bar V}, \frac{t}{2b}\Big)\Big).
\]
When all the $V_i$ are equal to $V$, $\bar V = V$.
\end{corollary}

\begin{proof}
\leanok
\uses{lem:subexp_sum, lem:subexp_const_mul, lem:subexp_tail}
Lemma~\ref{lem:subexp_sum} gives parameters $(\sum_iV_i, b) = (K\bar V, b)$ for the sum, and
Lemma~\ref{lem:subexp_const_mul} with $a = 1/K$ gives $(K\bar V/K^2, b/K) = (\bar V/K, b/K)$ for
the average. Lemma~\ref{lem:subexp_tail} (with $b/K > 0$ and $\bar V/K > 0$) then gives the
two-sided bound
$2\exp(-\min(t^2/(2\bar V/K), t/(2b/K))) = 2\exp(-K\min(t^2/(2\bar V), t/(2b)))$.
(For $b = 0$ the same computation with the $b = 0$ clause of Lemma~\ref{lem:subexp_tail} gives
$2\exp(-Kt^2/(2\bar V))$, Mathlib's \texttt{measure\_sum\_ge\_le\_of\_iIndepFun}.)
\end{proof}

\section{Squares of Gaussian variables}
\label{sec:pre_subexp_gaussian}

\begin{lemma}[MGF of the square of a real Gaussian]\label{lem:gaussian_square_mgf}
\lean{integral_exp_quadratic, integrable_exp_quadratic, ProbabilityTheory.integral_exp_mul_sq_gaussianReal, ProbabilityTheory.integrable_exp_mul_sq_gaussianReal}\leanok
\uses{def:pg_gaussian_vector}
For $b < 0$ and $c, d \in \R$,
\[
  \int_{\R} \exp(bx^2 + cx + d)\,dx = \sqrt{\frac{\pi}{-b}}\,\exp\Big(d - \frac{c^2}{4b}\Big),
\]
the integrand being integrable. Consequently, let $\mu \in \R$, $\sigma \ge 0$,
$X \sim \Normal(\mu,\sigma^2)$ and $\lambda \in \R$ with $2\sigma^2\lambda < 1$. Then
$e^{\lambda X^2}$ is integrable and
\[
  \E e^{\lambda X^2} = (1 - 2\sigma^2\lambda)^{-1/2}\exp\Big(\frac{\mu^2\lambda}{1 - 2\sigma^2\lambda}\Big).
\]
\end{lemma}

\begin{proof}
\leanok
\uses{def:pg_gaussian_vector}
Complete the square: $bx^2 + cx + d = b\,(x + \frac{c}{2b})^2 + d - \frac{c^2}{4b}$, so by the
translation invariance of Lebesgue measure (\texttt{integral\_add\_right\_eq\_self}) and
Mathlib \texttt{integral\_gaussian} ($\int e^{-\beta y^2}dy = \sqrt{\pi/\beta}$ with
$\beta = -b > 0$), the first identity follows; integrability is
\texttt{integrable\_exp\_neg\_mul\_sq} after the same translation.
If $\sigma = 0$ then $X = \mu$ a.s.\ and both sides of the second identity equal
$e^{\lambda\mu^2}$. Let $\sigma > 0$. With the density
$\frac{1}{\sqrt{2\pi\sigma^2}}e^{-(x-\mu)^2/(2\sigma^2)}$ (Mathlib \texttt{gaussianReal},
\texttt{gaussianPDFReal}, \texttt{integral\_gaussianReal\_eq\_integral\_smul}),
\[
  \E e^{\lambda X^2} = \frac{1}{\sqrt{2\pi\sigma^2}}\int_{\R}\exp\Big(\lambda x^2 - \frac{(x-\mu)^2}{2\sigma^2}\Big)dx ,
\]
and $\lambda x^2 - \frac{(x-\mu)^2}{2\sigma^2} = bx^2 + cx + d$ with
$b = \lambda - \frac{1}{2\sigma^2} = -\frac{1-2\sigma^2\lambda}{2\sigma^2} < 0$,
$c = \frac{\mu}{\sigma^2}$, $d = -\frac{\mu^2}{2\sigma^2}$. The first identity gives
$\frac{1}{\sqrt{2\pi\sigma^2}}\sqrt{\frac{2\pi\sigma^2}{1-2\sigma^2\lambda}}
= (1-2\sigma^2\lambda)^{-1/2}$ and the exponent
$d - \frac{c^2}{4b} = -\frac{\mu^2}{2\sigma^2} + \frac{\mu^2}{2\sigma^2(1-2\sigma^2\lambda)}
= \frac{\mu^2\lambda}{1-2\sigma^2\lambda}$.
\end{proof}

\begin{lemma}[Two elementary inequalities]\label{lem:log_ineq_nu}
\lean{ProbabilityTheory.neg_log_one_sub_div_two_sub_le_sq, ProbabilityTheory.one_div_one_sub_le_two}\leanok
For every $\nu \in \R$ with $\abs{\nu} \le 1/2$:
\begin{enumerate}
  \item[(i)] $-\frac12\log(1-\nu) - \frac{\nu}{2} \le \nu^2$;
  \item[(ii)] $\frac{1}{1-\nu} \le 2$.
\end{enumerate}
\end{lemma}

\begin{proof}
\leanok
(i) Mathlib's \texttt{Real.abs\_log\_sub\_add\_sum\_range\_le}, for $\abs x < 1$ and $n = 1$,
reads $\abs{x + \log(1-x)} \le \abs{x}^2/(1 - \abs x)$. With $x = \nu$ and $\abs\nu \le 1/2$,
$1 - \abs\nu \ge 1/2$, so $\abs{\nu + \log(1-\nu)} \le 2\nu^2$, hence
$-\log(1-\nu) - \nu \le 2\nu^2$ and dividing by $2$ gives (i).
(ii) $1 - \nu \ge 1 - \abs\nu \ge \frac12$.
\end{proof}

\begin{lemma}[Centered square of a Gaussian is sub-exponential]\label{lem:gaussian_square_subexp}
\lean{ProbabilityTheory.hasSubexponentialMGF_sq_sub_gaussianReal}\leanok
\uses{def:subexp, def:pg_gaussian_vector}
Let $\mu\in\R$, $\sigma \ge 0$ and $X \sim \Normal(\mu,\sigma^2)$. Then
\[
  \mathrm{HasSubexponentialMGF}\big(X^2 - (\mu^2 + \sigma^2),\ 8(\sigma^4 + \mu^2\sigma^2),\ 4\sigma^2\big),
\]
i.e.\ for all $\abs\lambda \le \frac{1}{4\sigma^2}$,
$\log\E e^{\lambda(X^2 - (\mu^2+\sigma^2))} \le 4(\sigma^4 + \mu^2\sigma^2)\lambda^2$
(and $\E X^2 = \mu^2 + \sigma^2$ by Lemma~\ref{lem:ps_subexp_integral_zero}).
\end{lemma}

\begin{proof}
\leanok
\uses{lem:gaussian_square_mgf, lem:log_ineq_nu}
Let $4\sigma^2\abs\lambda \le 1$ and $\nu := 2\sigma^2\lambda$, so that $\abs\nu \le 1/2$ and in
particular $2\sigma^2\lambda < 1$: Lemma~\ref{lem:gaussian_square_mgf} applies, $e^{\lambda X^2}$
is integrable (hence so is $e^{\lambda(X^2 - \mu^2 - \sigma^2)}$), and
\[
  \log\E e^{\lambda(X^2 - \mu^2 - \sigma^2)}
  = -\frac12\log(1-\nu) + \frac{\mu^2\lambda}{1-\nu} - \lambda(\mu^2 + \sigma^2)
  = \Big(-\frac12\log(1-\nu) - \frac\nu2\Big) + \mu^2\lambda\Big(\frac{1}{1-\nu} - 1\Big),
\]
using $\lambda\sigma^2 = \nu/2$. By Lemma~\ref{lem:log_ineq_nu}(i), the first bracket is at most
$\nu^2 = 4\sigma^4\lambda^2$. The second term equals
$\mu^2\lambda\nu/(1-\nu) = 2\mu^2\sigma^2\lambda^2\cdot\frac{1}{1-\nu} \le 4\mu^2\sigma^2\lambda^2$
by Lemma~\ref{lem:log_ineq_nu}(ii), since $\lambda\nu = 2\sigma^2\lambda^2 \ge 0$. Hence
$\log\E e^{\lambda(X^2-\mu^2-\sigma^2)} \le 4(\sigma^4 + \mu^2\sigma^2)\lambda^2 = \frac{\lambda^2}{2}\cdot8(\sigma^4 + \mu^2\sigma^2)$,
which is the claim with $V = 8(\sigma^4 + \mu^2\sigma^2)$ and $b = 4\sigma^2$.
\end{proof}

\begin{lemma}[Chi-square variables are sub-exponential]\label{lem:chi_square_subexp}
\lean{ProbabilityTheory.hasSubexponentialMGF_sum_sq_sub_one_pi_gaussianReal, ProbabilityTheory.hasSubexponentialMGF_norm_sq_sub_stdGaussian}\leanok
\uses{def:subexp, def:pg_gaussian_vector}
Let $g \sim \Normal(0, I_d)$. Then
$\mathrm{HasSubexponentialMGF}\big(\norm{g}^2 - d,\ 8d,\ 4\big)$.
\end{lemma}

\begin{proof}
\leanok
\uses{lem:gauss_pi, lem:gaussian_square_subexp, lem:subexp_sum, lem:ps_subexp_map}
$\norm{g}^2 - d = \sum_{i<d}(g_i^2 - 1)$ with $g_0,\dots,g_{d-1}$ i.i.d.\ $\Normal(0,1)$
(Lemma~\ref{lem:gauss_pi}; in Lean, $\Normal(0,I_d)$ is the image of the product measure
$\Normal(0,1)^{\otimes d}$ on $\R^d$, \texttt{map\_pi\_eq\_stdGaussian}, and the statement is
transported by Lemma~\ref{lem:ps_subexp_map}). Each $g_i^2 - 1$ has parameters $(8, 4)$ by
Lemma~\ref{lem:gaussian_square_subexp} with $\mu = 0$, $\sigma = 1$; the summands are
independent (\texttt{iIndepFun\_pi}), so Lemma~\ref{lem:subexp_sum} gives $(8d, 4)$.
\end{proof}

\begin{lemma}[Chi-square tail]\label{lem:chi_square_tail}
\lean{ProbabilityTheory.measureReal_norm_sq_sub_ge_le_stdGaussian, ProbabilityTheory.measureReal_norm_sq_ge_le_stdGaussian}\leanok
\uses{def:pg_gaussian_vector}
Let $g \sim \Normal(0,I_d)$ with $d \ge 1$. For every $t \ge 0$,
\[
  \Prob\big(\norm{g}^2 - d \ge t\big) \le \exp\Big(-\min\Big(\frac{t^2}{16d}, \frac{t}{8}\Big)\Big),
\]
and for every $\delta \in (0,1)$,
$\Prob\big(\norm{g}^2 \ge 2d + 12\log(1/\delta)\big) \le \delta$.
\end{lemma}

\begin{proof}
\leanok
\uses{lem:chi_square_subexp, lem:subexp_tail}
The first bound is Lemma~\ref{lem:subexp_tail} with $(V,b) = (8d,4)$ from
Lemma~\ref{lem:chi_square_subexp}: $t^2/(2V) = t^2/(16d)$ and $t/(2b) = t/8$. For the second,
let $L := \log(1/\delta) > 0$ and $t := \max(4\sqrt{dL}, 8L)$. Then $t^2/(16d) \ge 16dL/(16d) = L$
and $t/8 \ge L$, so $\Prob(\norm{g}^2 - d \ge t) \le e^{-L} = \delta$. Moreover, by AM--GM,
$4\sqrt{dL} = 2\sqrt{d\cdot4L} \le d + 4L$, so $t \le \max(d + 4L, 8L) \le d + 8L \le d + 12L$.
Hence $\{\norm{g}^2 \ge 2d + 12L\} = \{\norm{g}^2 - d \ge d + 12L\} \subseteq \{\norm{g}^2 - d \ge t\}$
has probability at most $\delta$.
\end{proof}

\begin{lemma}[Gaussian quadratic forms are sub-exponential]\label{lem:gauss_quadratic_form_subexp}
\uses{def:subexp, def:pg_gaussian_vector}
Let $S \in \PSD$ ($d\times d$, $S \ne 0$) and $g \sim \Normal(0,I_d)$. Then
$\E[g^\top Sg] = \tr S$ and
\[
  \mathrm{HasSubexponentialMGF}\big(g^\top Sg - \tr S,\ 8\tr(S^2),\ 4\norm{S}_{\mathrm{op}}\big),
\]
where $\norm{S}_{\mathrm{op}}$ is the largest eigenvalue of $S$.
\end{lemma}

\begin{proof}
\uses{lem:gauss_rotation, lem:gauss_pi, lem:gaussian_square_subexp, lem:subexp_const_mul, lem:ps_subexp_mono, lem:subexp_sum, lem:gauss_norm_moments}
Diagonalize $S = U\operatorname{diag}(s_1,\dots,s_d)U^\top$ with $U$ orthogonal and
$s_i \ge 0$ (Mathlib \texttt{Matrix.IsHermitian.spectral\_theorem}), so that
$\tr S = \sum_is_i$, $\tr(S^2) = \sum_is_i^2$ and $\norm{S}_{\mathrm{op}} = s_{\max} := \max_is_i > 0$.
Let $h := U^\top g$; then $h \sim \Normal(0,I_d)$ (Lemma~\ref{lem:gauss_rotation}(i)), its
coordinates are i.i.d.\ $\Normal(0,1)$ (Lemma~\ref{lem:gauss_pi}), and
$g^\top Sg - \tr S = \sum_is_i(h_i^2 - 1)$. The expectation is
Lemma~\ref{lem:gauss_norm_moments}(i). For each $i$ with $s_i > 0$,
Lemma~\ref{lem:gaussian_square_subexp} gives $(8,4)$ for $h_i^2 - 1$ and
Lemma~\ref{lem:subexp_const_mul} gives $(8s_i^2, 4s_i)$ for $s_i(h_i^2-1)$, hence
$(8s_i^2, 4s_{\max})$ by Lemma~\ref{lem:ps_subexp_mono}; for $s_i = 0$ the summand is $0$,
which has parameters $(0, 4s_{\max})$. The summands are independent, so
Lemma~\ref{lem:subexp_sum} gives $(8\sum_is_i^2, 4s_{\max}) = (8\tr(S^2), 4\norm{S}_{\mathrm{op}})$.
\end{proof}

\begin{remark}[Not used]\label{rem:ps_quadratic_form_unused}
Lemma~\ref{lem:gauss_quadratic_form_subexp} is recorded for completeness (it is the general
form of Lemma~\ref{lem:chi_square_subexp}, which is the case $S = I_d$) but none of the main
results uses it: the large-norm regime of the norm estimation uses a deterministic basis
design, for which only the chi-square case is needed. It is not formalized.
\end{remark}

\section{Squares of sub-Gaussian variables and of Rademacher linear forms}
\label{sec:pre_subexp_rademacher}

\begin{lemma}[MGF of the square of a sub-Gaussian variable, positive parameter]
\label{lem:subgaussian_square_mgf_pos}
\lean{ProbabilityTheory.HasSubgaussianMGF.lintegral_exp_mul_sq_le, ProbabilityTheory.HasSubgaussianMGF.integrable_exp_mul_sq, ProbabilityTheory.HasSubgaussianMGF.integral_exp_mul_sq_le}\leanok
\uses{def:pg_gaussian_vector}
Let $Z$ be a real random variable with a sub-Gaussian MGF with constant $c \ge 0$
(Mathlib \texttt{HasSubgaussianMGF Z c}), and let $0 \le \lambda$ with $2c\lambda < 1$
(any $\lambda \ge 0$ if $c = 0$). Then $e^{\lambda Z^2}$ is integrable and
\[
  \E e^{\lambda Z^2} \le (1 - 2c\lambda)^{-1/2}.
\]
\end{lemma}

\begin{proof}
\leanok
\uses{lem:gaussian_square_mgf, lem:gauss_1d_moments}
\emph{Gaussian trick:} for every real $z$, by the MGF of $g' \sim \Normal(0,1)$ at
$s = \sqrt{2\lambda}\,z$ (Mathlib \texttt{mgf\_gaussianReal}),
\[
  e^{\lambda z^2} = e^{(\sqrt{2\lambda}\,z)^2/2} = \E_{g'}\exp\big(\sqrt{2\lambda}\,z\,g'\big).
\]
Hence, by Tonelli's theorem for the nonnegative measurable function
$(\omega, v) \mapsto \exp(\sqrt{2\lambda}\,Z(\omega)\,v)$ on the product of $(\Omega,\Prob)$
with $(\R, \Normal(0,1))$ (\texttt{MeasureTheory.lintegral\_lintegral\_swap}),
\[
  \E e^{\lambda Z^2} = \int\Big(\int\exp\big(\sqrt{2\lambda}\,Z(\omega)\,v\big)\,\Normal(0,1)(dv)\Big)\Prob(d\omega)
  = \int\Big(\int\exp\big((\sqrt{2\lambda}\,v)\,Z(\omega)\big)\Prob(d\omega)\Big)\Normal(0,1)(dv).
\]
For fixed $v$, the inner integral is the MGF of $Z$ at $s := \sqrt{2\lambda}\,v$, which is at
most $e^{cs^2/2} = e^{c\lambda v^2}$ by sub-Gaussianity. Therefore
\[
  \E e^{\lambda Z^2} \le \int e^{c\lambda v^2}\,\Normal(0,1)(dv) = (1 - 2c\lambda)^{-1/2}
\]
by Lemma~\ref{lem:gaussian_square_mgf} with $\mu = 0$, $\sigma = 1$ and the parameter
$c\lambda < 1/2$. Since the right-hand side is finite, $e^{\lambda Z^2}$ is integrable: all
integrals above are Lebesgue integrals in $[0,\infty]$ of nonnegative functions, and
\texttt{Integrable} is the conjunction of \texttt{AEStronglyMeasurable} (clear) and
\texttt{HasFiniteIntegral}, which for a nonnegative function $f$ is
$\int^-\mathrm{ENNReal.ofReal}(f)\,d\Prob < \infty$ (Mathlib \texttt{hasFiniteIntegral\_iff\_ofReal}).
\end{proof}

\begin{lemma}[Fourth moment of a sub-Gaussian variable]\label{lem:subgaussian_fourth_moment}
\lean{ProbabilityTheory.HasSubgaussianMGF.integral_pow_four_le}\leanok
Let $Z$ have a sub-Gaussian MGF with constant $c \ge 0$ under a probability measure. Then
$\E Z^4 \le 14c^2$.
\end{lemma}

\begin{proof}
\leanok
\uses{lem:subgaussian_square_mgf_pos}
If $c = 0$ then $Z = 0$ a.s.\ (Mathlib \texttt{ae\_eq\_zero\_of\_hasSubgaussianMGF\_zero}). Let
$c > 0$ and $\lambda := 1/(4c)$, so that $2c\lambda = 1/2$ and
Lemma~\ref{lem:subgaussian_square_mgf_pos} gives $\E e^{\lambda Z^2} \le \sqrt2$. For $x \ge 0$,
$1 + x + x^2/2 \le e^x$ (Mathlib \texttt{Real.quadratic\_le\_exp\_of\_nonneg}), so with
$x = \lambda Z^2 \ge 0$, pointwise $\frac{\lambda^2}{2}Z^4 \le e^{\lambda Z^2} - 1$. Taking
expectations, $\frac{\lambda^2}{2}\E Z^4 \le \sqrt2 - 1$, i.e.\
$\E Z^4 \le 32c^2(\sqrt2 - 1) \le 14c^2$ since $\sqrt2 \le 23/16$.
\end{proof}

\begin{lemma}[A quadratic upper bound for the exponential on the negative half-line]
\label{lem:ps_exp_le_quadratic}
\lean{Real.exp_le_one_add_add_sq_div_two_of_nonpos}\leanok
For every $x \le 0$, $e^x \le 1 + x + \frac{x^2}{2}$.
\end{lemma}

\begin{proof}
\leanok
Let $y := -x \ge 0$. By Mathlib \texttt{Real.quadratic\_le\_exp\_of\_nonneg},
$e^{y} \ge 1 + y + y^2/2 > 0$, so $e^x = 1/e^y \le 1/(1 + y + y^2/2)$. Finally
$1/(1 + y + y^2/2) \le 1 - y + y^2/2$, because
$(1 + y + y^2/2)(1 - y + y^2/2) = (1 + y^2/2)^2 - y^2 = 1 + y^4/4 \ge 1$ and
$1 - y + y^2/2 = ((1-y)^2 + 1)/2 > 0$. Since $1 - y + y^2/2 = 1 + x + x^2/2$, this is the claim.
\end{proof}

\begin{lemma}[Centered square of a sub-Gaussian variable is sub-exponential]
\label{lem:subgaussian_square_subexp}
\lean{ProbabilityTheory.HasSubgaussianMGF.hasSubexponentialMGF_sq_sub}\leanok
\uses{def:subexp}
Let $Z$ have a sub-Gaussian MGF with constant $c \ge 0$ under a probability measure, and
assume $\E Z^2 = c$. Then
\[
  \mathrm{HasSubexponentialMGF}\big(Z^2 - c,\ 16c^2,\ 4c\big).
\]
\end{lemma}

\begin{proof}
\leanok
\uses{lem:subgaussian_square_mgf_pos, lem:subgaussian_fourth_moment, lem:log_ineq_nu, lem:ps_exp_le_quadratic}
$Z$ has all moments (Mathlib \texttt{integrable\_pow\_of\_mem\_interior\_integrableExpSet}).
Let $\abs\lambda \le 1/(4c)$; we show that $e^{\lambda(Z^2 - c)}$ is integrable and
$\E e^{\lambda(Z^2 - c)} \le e^{8c^2\lambda^2} = e^{\lambda^2\cdot16c^2/2}$.

\emph{Case $0 \le \lambda \le 1/(4c)$.} Let $\nu := 2c\lambda \in [0, 1/2]$, so that
$2c\lambda < 1$ and Lemma~\ref{lem:subgaussian_square_mgf_pos} gives the integrability of
$e^{\lambda Z^2}$ (hence of $e^{\lambda(Z^2-c)} = e^{-\lambda c}e^{\lambda Z^2}$) and
$\E e^{\lambda Z^2} \le (1-\nu)^{-1/2}$. Hence, using $\lambda c = \nu/2$ and
Lemma~\ref{lem:log_ineq_nu}(i),
\[
  \log\E e^{\lambda(Z^2 - c)} \le -\frac12\log(1-\nu) - \frac{\nu}{2} \le \nu^2 = 4c^2\lambda^2 \le 8c^2\lambda^2 .
\]

\emph{Case $-1/(4c) \le \lambda < 0$.} Pointwise $\lambda Z^2 \le 0$, so
$e^{\lambda(Z^2-c)} \le e^{-\lambda c}$ is bounded, hence integrable, and
Lemma~\ref{lem:ps_exp_le_quadratic} gives $e^{\lambda Z^2} \le 1 + \lambda Z^2 + \frac{\lambda^2Z^4}{2}$.
Taking expectations and using $\E Z^2 = c$ and Lemma~\ref{lem:subgaussian_fourth_moment}
($\E Z^4 \le 14c^2$),
\[
  \E e^{\lambda Z^2} \le 1 + \lambda c + 7c^2\lambda^2
  \le \exp\big(\lambda c + 7c^2\lambda^2\big),
\]
by $1 + y \le e^y$. Multiplying by $e^{-\lambda c}$:
$\E e^{\lambda(Z^2 - c)} \le e^{7c^2\lambda^2} \le e^{8c^2\lambda^2}$.

In both cases $\E e^{\lambda(Z^2-c)} \le \exp(\lambda^2\cdot16c^2/2)$, which is the claim with
$V = 16c^2$ and $b = 4c$.
\end{proof}

\begin{definition}[Rademacher vector]\label{def:ps_rademacher}
\lean{ProbabilityTheory.signMeasure, ProbabilityTheory.rademacherMeasure, ProbabilityTheory.integral_signMeasure, ProbabilityTheory.iIndepFun_eval_rademacherMeasure, ProbabilityTheory.map_eval_rademacherMeasure, ProbabilityTheory.ae_rademacherMeasure_mem_Icc, ProbabilityTheory.integral_eval_rademacherMeasure, ProbabilityTheory.integral_eval_sq_rademacherMeasure, ProbabilityTheory.inner_toLp_eq_sum}\leanok
A \emph{Rademacher vector} in $\R^d$ is a random vector $\eps = (\eps_1,\dots,\eps_d)$ whose
coordinates are independent with $\Prob(\eps_i = 1) = \Prob(\eps_i = -1) = 1/2$; its law is the
product $\big(\frac12\delta_{1} + \frac12\delta_{-1}\big)^{\otimes d}$
(Lean: \texttt{signMeasure := 2⁻¹ • (dirac 1 + dirac (-1))} and
\texttt{rademacherMeasure ι := Measure.pi (fun \_ ↦ signMeasure)} on \texttt{ι → ℝ}; the
coordinates are independent by \texttt{iIndepFun\_pi}, each with law \texttt{signMeasure},
$\E\eps_i = 0$, $\E\eps_i^2 = 1$ and $\abs{\eps_i} \le 1$ a.s.).
For $\theta \in \R^d$ we consider the linear form $Z := \ip{\eps}{\theta} = \sum_i\eps_i\theta_i$,
which satisfies $\abs Z \le \sum_i\abs{\theta_i}$.
\end{definition}

\begin{lemma}[Rademacher linear forms are sub-Gaussian]\label{lem:rademacher_linear_subgaussian}
\lean{ProbabilityTheory.hasSubgaussianMGF_sum_mul_rademacherMeasure, ProbabilityTheory.hasSubgaussianMGF_sum_mul_rademacherMeasure'}\leanok
\uses{def:ps_rademacher}
Let $\eps$ be a Rademacher vector in $\R^d$ and $\theta \in \R^d$. Then $\ip{\eps}{\theta}$ has a
sub-Gaussian MGF with constant $\norm{\theta}^2$ (Mathlib
\texttt{HasSubgaussianMGF ⟪ε, θ⟫ ‖θ‖²}): $\E e^{\lambda\ip{\eps}{\theta}} \le e^{\norm\theta^2\lambda^2/2}$
for all $\lambda\in\R$.
\end{lemma}

\begin{proof}
\leanok
\uses{def:ps_rademacher}
Each $\theta_i\eps_i$ takes values in $[-\abs{\theta_i}, \abs{\theta_i}]$ and has mean zero, so
by Hoeffding's lemma (Mathlib \texttt{hasSubgaussianMGF\_of\_mem\_Icc\_of\_integral\_eq\_zero})
it has a sub-Gaussian MGF with constant $((2\abs{\theta_i})/2)^2 = \theta_i^2$. The
$\theta_i\eps_i$ are independent, so \texttt{HasSubgaussianMGF.sum\_of\_iIndepFun} gives the
constant $\sum_i\theta_i^2 = \norm\theta^2$ for the sum.
\end{proof}

\begin{lemma}[First and second moments of a Rademacher linear form]\label{lem:rademacher_fourth_moment}
\lean{ProbabilityTheory.integral_eval_mul_eval_rademacherMeasure, ProbabilityTheory.integral_sum_mul_rademacherMeasure, ProbabilityTheory.integral_sq_sum_mul_rademacherMeasure}\leanok
\uses{def:ps_rademacher}
Let $\eps$ be a Rademacher vector in $\R^d$, $\theta\in\R^d$ and $Z := \ip{\eps}{\theta}$. Then
$\E[\eps_i\eps_j] = \indic\{i = j\}$, $\E Z = 0$ and $\E Z^2 = \norm{\theta}^2$.
\end{lemma}

\begin{proof}
\leanok
\uses{def:ps_rademacher}
All the variables involved are bounded, hence integrable. $\E[\eps_i^2] = 1$ and, for
$i \ne j$, $\E[\eps_i\eps_j] = \E\eps_i\,\E\eps_j = 0$ by independence. Then
$\E Z = \sum_i\theta_i\E\eps_i = 0$ and, expanding the square,
$\E Z^2 = \sum_{i,j}\theta_i\theta_j\E[\eps_i\eps_j] = \sum_i\theta_i^2$.
\end{proof}

\begin{lemma}[Centered square of a Rademacher linear form is sub-exponential]
\label{lem:rademacher_square_subexp}
\lean{ProbabilityTheory.hasSubexponentialMGF_sq_sum_mul_sub_rademacherMeasure}\leanok
\uses{def:subexp, def:ps_rademacher}
Let $\eps$ be a Rademacher vector in $\R^d$, $\theta \in \R^d$ with $r := \norm{\theta}$, and
$Z := \ip{\eps}{\theta}$. Then
\[
  \mathrm{HasSubexponentialMGF}\big(Z^2 - r^2,\ 16r^4,\ 4r^2\big).
\]
\end{lemma}

\begin{proof}
\leanok
\uses{lem:rademacher_linear_subgaussian, lem:rademacher_fourth_moment, lem:subgaussian_square_subexp}
By Lemma~\ref{lem:rademacher_linear_subgaussian}, $Z$ is sub-Gaussian with constant
$c = r^2$, and $\E Z^2 = r^2 = c$ by Lemma~\ref{lem:rademacher_fourth_moment}. This is exactly
the hypothesis of Lemma~\ref{lem:subgaussian_square_subexp}, which gives the parameters
$(16c^2, 4c) = (16r^4, 4r^2)$.
\end{proof}

\begin{remark}[Constants, and the case $r = 0$]\label{rem:ps_rademacher_constants}
The proof of Lemma~\ref{lem:subgaussian_square_subexp} shows the stronger parameters
$(8c^2, 4c)$ (the bound $e^{7c^2\lambda^2}$ for negative $\lambda$ could be improved by using
the exact fourth moment $\E Z^4 = 3r^4 - 2\sum_i\theta_i^4$ of a Rademacher form instead of
the general bound $14c^2$, but this is not needed); the value $V = 16r^4$ is kept because it is
the one used by Chapter~\ref{chap:norm} (Lemma~\ref{lem:term1_subexp}). The paper obtains
parameters $(Cr^4, Cr^2)$ with an unspecified $C$ from the Hanson--Wright inequality (its
Lemma \emph{lem:hanson-wright}, Appendix~A) and a tail-to-MGF conversion; the argument above
is self-contained and explicit. If $r = 0$ then $Z = 0$ a.s.\ and $Z^2 - r^2 = 0$, which has
parameters $(0, 0)$, consistent with the statement; the norm-estimation chapter treats this
case separately anyway.
\end{remark}

\textbf{Lean remark.} The Gaussian trick of Lemma~\ref{lem:subgaussian_square_mgf_pos} needs an
auxiliary standard Gaussian independent of $Z$; it is encoded by Tonelli's theorem on the
product measure \texttt{μ.prod (gaussianReal 0 1)}, without introducing a new random variable.
The Rademacher results are stated for the product measure \texttt{rademacherMeasure ι} on
\texttt{ι → ℝ} and the coordinate sum \texttt{∑ i, ε i * θ i}; the identification with the
inner product on \texttt{EuclideanSpace ℝ ι} is \texttt{inner\_toLp\_eq\_sum}.
